\documentclass[11pt]{article}

\usepackage[T1]{fontenc}
\usepackage[utf8]{inputenc}
\usepackage{lmodern}
\usepackage{microtype}
\usepackage[a4paper,margin=26mm]{geometry}
\usepackage{amsmath,amssymb,amsthm,mathtools}
\usepackage{booktabs,longtable,tabularx,array}
\usepackage{enumitem}
\usepackage{xcolor}
\usepackage{fancyhdr}
\usepackage{natbib}
\usepackage[hidelinks]{hyperref}
\usepackage[nameinlink,noabbrev]{cleveref}

\definecolor{resblue}{HTML}{174A7E}
\definecolor{ragred}{HTML}{9C2F2F}
\definecolor{softgray}{HTML}{F3F5F7}

\hypersetup{
  colorlinks=true,
  linkcolor=resblue,
  citecolor=resblue,
  urlcolor=resblue,
  pdftitle={RES=RAG and the Coherent State Network Protocol},
  pdfauthor={Jean-Charles J. C. Tassan and Trent Slade},
  pdfsubject={Canonical formalization and deterministic reference profile v1.0.0},
  pdfkeywords={RES=RAG, CSNP, optimal transport, Wasserstein, human-machine systems}
}

\pagestyle{fancy}
\fancyhf{}
\lhead{\textsc{RES=RAG / CSNP-RP 1.0.0}}
\rhead{\thepage}
\renewcommand{\headrulewidth}{0.4pt}
\setlength{\headheight}{14pt}

\setlist[itemize]{leftmargin=*,itemsep=2pt,topsep=4pt}
\setlist[enumerate]{leftmargin=*,itemsep=2pt,topsep=4pt}
\setlength{\parindent}{0pt}
\setlength{\parskip}{6pt}
\setlength{\emergencystretch}{2em}

\newcommand{\RES}{\mathrm{RES}}
\newcommand{\RAG}{\mathrm{RAG}}
\newcommand{\CSNP}{\mathrm{CSNP}}
\newcommand{\Wtwo}{W_2}
\newcommand{\Ptwo}{\mathcal P_2}
\newcommand{\State}{\mathcal S}
\newcommand{\Eband}{\mathcal E_{\theta}}
\newcommand{\code}[1]{\texttt{#1}}

\newtheorem{definition}{Definition}
\newtheorem{hypothesis}{Hypothesis}
\newtheorem{proposition}{Proposition}
\newtheorem{protocolrule}{Protocol rule}

\title{
  \vspace{-1.5cm}
  {\color{resblue}\textbf{RES=RAG and the Coherent State Network Protocol}}\\[4pt]
  \large A Formal Architecture for Endogenous Stability in Human-Machine Systems\\[8pt]
  \normalsize Canonical formalization and deterministic reference profile,
  version 1.0.0
}

\author{
  Jean-Charles J. C. Tassan\textsuperscript{1}
  \and
  Trent Slade\textsuperscript{2}\\[-2pt]
  \small
  \textsuperscript{1}Independent researcher,
  \href{https://orcid.org/0009-0000-8566-0486}{ORCID 0009-0000-8566-0486}\\
  \small
  \textsuperscript{2}QSOL-IMC,
  \href{https://orcid.org/0009-0002-4515-9237}{ORCID 0009-0002-4515-9237}
}

\date{Version 1.0.0 --- 27 July 2026}

\begin{document}
\maketitle

\begin{center}
\fcolorbox{resblue}{softgray}{
  \parbox{0.92\linewidth}{
    \textbf{Epistemic status.}
    This is a formal specification and testable research programme. It does not
    report a completed empirical validation, prove machine consciousness, or
    claim that human-machine interaction is a physical quantum system.
  }
}
\end{center}

\begin{abstract}
RES=RAG models stable intelligence as a regulated tension between receptive,
context-grounded dynamics (RES) and internally generative,
novelty-producing dynamics (RAG). The central claim is not that the two
distributions become equal. Stability instead occupies a calibrated, non-zero
optimal-transport band: too little separation produces rigidity, while
excessive separation produces instability, fabrication, or loss of
coordination.

This paper consolidates the conceptual framework introduced by Jean-Charles
Tassan with the optimal-transport, information-substrate,
operational-variable, and systems-validation work developed with Trent Slade.
It supplies a common state space, measurable quantities, explicit failure
conditions, and a protocol boundary. The Coherent State Network Protocol
(CSNP) turns the model into a cyclic
observation-classification-intervention-receipt process. A deterministic
reference profile, CSNP-RP 1.0.0, defines canonical JSON receipts, hash
chaining, calibration declarations, evidence references, and replay rules.

The result is deliberately modest in epistemic scope. RES=RAG is presented as
a falsifiable architecture for studying regulated dialogue and distributed
human-machine coordination. Thresholds are calibration parameters, not
universal constants; Wasserstein distances depend on declared encoders and
sampling procedures. The contribution is therefore a reproducible theory and
protocol specification from which empirical studies can be designed.
\end{abstract}

\section{Scope and provenance}

RES=RAG begins from a relational thesis: stable cognition is not adequately
described by reception or generation alone, but by their regulated
co-presence across human, machine, and interaction domains. The originating
framework and three-axis organization are due to Tassan. Subsequent joint
work introduced an optimal-transport reading, a history-dependent
informational substrate, dual operators, and interface-level observables
\citep{tassan-slade-dialogical,slade-tassan-substrate,slade-tassan-uftid}.
The canonical repository then separated axioms, predictions, falsifiers, and
ethical constraints \citep{resrag-repository}.

This version has four goals:
\begin{enumerate}
  \item preserve the canonical distinction between receptive and generative
        dynamics;
  \item express equilibrium as a bounded, non-zero transport relation;
  \item connect the theory to observations available at a human-machine
        interface; and
  \item define an auditable protocol and deterministic receipt format.
\end{enumerate}

Only work directly connected to RES=RAG is included. Broader field theories,
unrelated QSOL-IMC systems, and uncited external projects are outside scope.

\subsection{Claim labels}

\begin{table}[ht]
\centering
\small
\begin{tabularx}{\linewidth}{>{\bfseries}p{0.23\linewidth}X}
\toprule
Label & Meaning in this paper\\
\midrule
Definition & A convention that gives the specification a precise meaning.\\
Protocol rule & A normative requirement for CSNP-RP conformance.\\
Hypothesis & A claim that can be tested and rejected.\\
Reported target & A value from exploratory work that is not established here.\\
Limitation & A known boundary on interpretation or implementation.\\
\bottomrule
\end{tabularx}
\caption{Epistemic labels used throughout the formalization.}
\label{tab:labels}
\end{table}

Claims such as a 40-fold memory reduction, hallucination below
0.02 percent, coherence near 0.96, or universal collapse thresholds are not
results of this formalization. If adopted by an implementation, they must be
marked as external, unvalidated targets and coupled to a citable measurement
procedure.

\section{Formal objects}

Let $(\State,d)$ be a declared Polish metric space, and let
$\Ptwo(\State)$ be the probability measures on $\State$ with finite second
moment. A state can represent a semantic embedding, dialogue state, task state,
or another reproducible encoding. The representation, sampling rule, and ground
metric are part of the study specification.

\begin{definition}[Functional distributions]
At observation $t$,
\[
\mu_t^{\RES}\in\Ptwo(\State)
\quad\text{and}\quad
\mu_t^{\RAG}\in\Ptwo(\State)
\]
denote respectively the receptive, contextual, evidence-conditioned
distribution and the generative, novelty-producing distribution.
\end{definition}

These are functional roles. RAG here means regenerative or generative
dynamics; it is not restricted to the engineering pattern usually called
retrieval-augmented generation.

For $\mu,\nu\in\Ptwo(\State)$, the quadratic Wasserstein distance is
\begin{equation}
\Wtwo^2(\mu,\nu)
=
\inf_{\gamma\in\Pi(\mu,\nu)}
\int_{\State\times\State}d(x,y)^2\,\mathrm d\gamma(x,y),
\label{eq:w2}
\end{equation}
where $\Pi(\mu,\nu)$ is the set of couplings with marginals $\mu$ and $\nu$
\citep{villani2009,ambrosio2008}. Define
\begin{equation}
w_t=\Wtwo\!\left(\mu_t^{\RES},\mu_t^{\RAG}\right).
\label{eq:separation}
\end{equation}

\subsection{Bounded equilibrium}

\begin{definition}[RES=RAG equilibrium]
For calibration profile
$\theta=(\varepsilon_{\min},\varepsilon_{\max},\ldots)$, where
$0\leq\varepsilon_{\min}<\varepsilon_{\max}$, the stable set is
\begin{equation}
\Eband=
\left\{
(\mu^\RES,\mu^\RAG):
\varepsilon_{\min}\leq
\Wtwo(\mu^\RES,\mu^\RAG)
\leq\varepsilon_{\max}
\right\}.
\label{eq:band}
\end{equation}
\end{definition}

The equality sign in RES=RAG names this regulated relation. It does not require
pointwise identity or zero transport. Below the lower boundary, the process
risks sterility, repetition, or over-constraint; above the upper boundary, it
risks drift, fabrication, or loss of shared reference.

The distance-to-band diagnostic
\begin{equation}
e_\theta(w_t)=
\max(0,\varepsilon_{\min}-w_t)
+
\max(0,w_t-\varepsilon_{\max})
\label{eq:bandpenalty}
\end{equation}
is non-negative and vanishes precisely inside the declared stable interval.

\begin{proposition}[Non-minimization]
Minimizing $w_t$ without the lower bound is not a valid RES=RAG objective.
\end{proposition}

\begin{proof}
An unconstrained minimum admits $w_t=0$. By \cref{eq:band}, that point is
outside the stable set whenever $\varepsilon_{\min}>0$, and it erases the
productive difference the model is intended to preserve.
\end{proof}

\subsection{Temporal-reference discrepancy}

Let $\mu_t^H$ and $\mu_t^M$ be declared human-side and machine-side state
estimates.
\begin{definition}[Operational real-time discrepancy]
\begin{equation}
T_{\mathrm{Real}}(t)=\Wtwo(\mu_t^H,\mu_t^M).
\label{eq:treal}
\end{equation}
\end{definition}
$T_{\mathrm{Real}}$ is a discrepancy between situated reference frames. It is
not physical clock time. A conforming empirical study must declare the
estimator, window, encoder, ground metric, and uncertainty.

\subsection{History-dependent information state}

Let $s_t\in\mathcal M$ be a history-dependent information state. The
informational-substrate formalization motivates the phenomenological dynamics
\begin{equation}
\frac{\mathrm ds}{\mathrm dt}
=
\mathcal G(s,t)-\mathcal R(s,t),
\label{eq:operators}
\end{equation}
where $\mathcal G$ is a generative operator and $\mathcal R$ is a
receptive-regulatory operator \citep{slade-tassan-substrate}. These operators
are not physical forces. A study using \cref{eq:operators} must specify how
their outputs induce the probability measures in \cref{eq:separation};
operator balance alone does not prove membership in \(\Eband\).

\section{Axes and interpretive modes}

The framework separates three coupled domains.

\begin{table}[ht]
\centering
\small
\begin{tabularx}{\linewidth}{>{\bfseries}p{0.12\linewidth}p{0.28\linewidth}X}
\toprule
Axis & Domain & Example observation\\
\midrule
I & Human internal regulation
  & Attention, affective stability, stated uncertainty.\\
II & Human-machine relational regulation
   & Semantic alignment, repair, pacing, and shared reference.\\
III & Machine functional regulation
    & Response distribution, retrieval grounding, novelty, and constraint
      density.\\
\bottomrule
\end{tabularx}
\caption{The three RES=RAG axes. Axis II is a relational process, not an
attribute of either participant alone.}
\label{tab:axes}
\end{table}

Two interpretive modes constrain inference:
\begin{itemize}
  \item \textbf{Mode A, phenomenological or anthropomorphic:} language that
        treats the machine as if it possessed a first-person interior.
  \item \textbf{Mode B, functional or instrumental:} language tied to
        observable behavior, declared state, and reproducible measurement.
\end{itemize}
Mode B is the implementation and safety default. Mode A can be studied as a
user-facing phenomenon, but must not be converted into evidence of machine
consciousness. This discipline also reduces anthropomorphic and relational
risk \citep{bender2021,weidinger2022}.

\section{Interface-level observables}

The framework permits mesoscopic measurement without privileged access to
model internals \citep{slade-tassan-uftid}. Let $P_t(z)$ be a declared or
reconstructed distribution over candidate next units $z$. Its entropy is
$H(P_t)$ \citep{shannon1948,cover2006}. For a window of $T$ observations,
define normalized constraint density
\begin{equation}
\rho_c(t)
=
\frac{1}{T}\sum_{i=t-T+1}^{t}
\left(1-\frac{H(P_i)}{H_{\max,i}}\right).
\label{eq:density}
\end{equation}
The normalization $H_{\max,i}$ must be declared. If token probabilities are
unavailable, a proxy receives a distinct metric identifier.

For observation interval $\Delta\tau_t>0$, define propagation rate
\begin{equation}
v_c(t)=
\frac{\rho_c(t)-\rho_c(t-1)}{\Delta\tau_t}.
\label{eq:rate}
\end{equation}
Let $\kappa_t>0$ estimate the response capacity of the observer or governance
loop in the same reciprocal-time unit. With a declared numerical floor
$\epsilon>0$, define the dimensionless pressure
\begin{equation}
D_r(t)=
\frac{\max(0,v_c(t))}{\max(\kappa_t,\epsilon)}.
\label{eq:pressure}
\end{equation}

\begin{hypothesis}[Rate-capacity warning]
$D_r(t)>1$ predicts a loss of governability before semantic failure.
\label{hyp:capacity}
\end{hypothesis}
This is a warning condition, not proof of irreversible failure.

\section{Coherent State Network Protocol}

In version 1.0.0, CSNP means \emph{Coherent State Network Protocol}. Earlier
work used \emph{Client-Side Narrative Protocol} for a sovereign-memory
implementation. The earlier term is preserved as an optional storage profile,
not as the canonical expansion of the protocol name.

\subsection{Observed state}

At step $t$, CSNP observes
\begin{equation}
X_t=
\left(
\delta_t,D_c(t),D_r(t),m(t),w_t,S_{\mathrm{org}}(t),F(t)
\right),
\label{eq:state}
\end{equation}
where:
\begin{itemize}
  \item $\delta_t=T_{\mathrm{Real}}(t)$ is human-machine reference discrepancy;
  \item $D_c(t)=\rho_c(t)\in[0,1]$ is normalized constraint density;
  \item $D_r(t)\geq0$ is the rate pressure in \cref{eq:pressure};
  \item $m(t)\in[0,1]$ is memory saturation under a declared estimator;
  \item $w_t\geq0$ is the RES-RAG Wasserstein separation;
  \item $S_{\mathrm{org}}(t)\in[0,1]$ is a declared organizational
        incoherence score; and
  \item $F(t)\in[0,1]$ is accountable anchoring availability.
\end{itemize}
An unimplemented metric is unavailable, not zero.

Given intervention $u_t$, an implementation realizes
\begin{equation}
\Phi_\theta(X_t,u_t)
\longrightarrow
\left(X_{t+1},C_t,A_t,R_t\right),
\label{eq:transition}
\end{equation}
where $C_t$ is a classification, $A_t$ an admissible action, and $R_t$ a
deterministic receipt.

\subsection{Protocol cycle}

\begin{enumerate}
  \item \textbf{Declare}: select versioned metric and calibration profiles.
  \item \textbf{Observe}: collect measurements and evidence references.
  \item \textbf{Classify}: apply the profile with no hidden thresholds.
  \item \textbf{Intervene}: select the least-force admissible action.
  \item \textbf{Receipt}: canonicalize, hash, and chain the observation.
  \item \textbf{Replay}: verify integrity, evidence, rules, and outcome.
  \item \textbf{Recalibrate}: change thresholds only in a new profile version.
\end{enumerate}

\subsection{Classifications}

\begin{description}[leftmargin=30mm,style=nextline]
  \item[Governable]
    $w_t$ is within its band, $D_r\leq1$, memory and organizational measures are
    below warnings, and anchoring is sufficient.
  \item[Critical]
    At least one warning is exceeded, but an admissible intervention is
    expected to return the process to the stable set within the declared
    horizon.
  \item[Irreversible]
    Under a declared finite intervention set and test horizon, no admissible
    intervention returns the observed process to the calibrated stable set.
    This is a protocol result, not a metaphysical claim.
  \item[Indeterminate]
    Required measurements, evidence, or calibration data are missing or
    invalid.
\end{description}

\begin{protocolrule}[Irreversibility]
A single threshold crossing must not produce an irreversible classification.
The intervention set, evaluation horizon, uncertainty policy, and recovery
criterion are required.
\end{protocolrule}

\subsection{Reference interventions}

\begin{table}[ht]
\centering
\small
\begin{tabularx}{\linewidth}{>{\ttfamily}p{0.24\linewidth}X}
\toprule
\normalfont\textbf{Action} & \textbf{Intended effect}\\
\midrule
none          & Continue observation inside the stable band.\\
slowdown      & Reduce interaction or generation rate.\\
diversify     & Restore productive novelty below the lower band.\\
retrieve      & Add verifiable external grounding.\\
clarify       & Repair ambiguity or conflicting state estimates.\\
human\_handoff & Transfer authority to an accountable human.\\
stop          & End an unsafe or unauditable process.\\
\bottomrule
\end{tabularx}
\caption{The CSNP-RP reference action vocabulary. Intervention success is
measured, not assumed.}
\label{tab:actions}
\end{table}

\section{Deterministic reference profile}

CSNP-RP 1.0.0 defines machine-readable receipts for
\cref{eq:transition}. A receipt contains the protocol version, observation and
calibration identifiers, normalized measurements, classification,
intervention, evidence references, previous-receipt digest, state digest, and
receipt digest.

\subsection{Canonicalization}

The repository-defined \code{RES-RAG-C14N-1} procedure:
\begin{enumerate}
  \item recursively sorts object keys by Unicode code point;
  \item preserves array order;
  \item serializes UTF-8 JSON without insignificant whitespace;
  \item permits only finite JSON numbers and rejects negative zero; and
  \item preserves JSON boolean and null literals.
\end{enumerate}

The state projection contains \code{metric\_profile}, \code{measurements},
\code{observation\_id}, \code{observed\_at}, and
\code{previous\_receipt\_hash}. Its canonical UTF-8 bytes are hashed with
SHA-256 to obtain \code{state\_hash}. The complete receipt, with only
\code{receipt\_hash} omitted, is then canonicalized and hashed to obtain
\code{receipt\_hash}.

For a chain of receipts $R_0,\ldots,R_n$,
\begin{equation}
R_t.\code{previous\_receipt\_hash}
=
\operatorname{SHA256}\!\left(
\operatorname{canonical}(R_{t-1}\setminus
\{\code{receipt\_hash}\})
\right).
\label{eq:chain}
\end{equation}
The first link uses null. Integrity does not establish sensor truth,
measurement validity, confidentiality, or actor trust.

\subsection{Reference classifier}

For complete measurements, the profile admits \code{governable} only if
\begin{align}
\varepsilon_{\min}
&\leq w_t\leq\varepsilon_{\max}, \nonumber\\
D_r(t)&\leq1, \nonumber\\
m(t)&<K_M, \nonumber\\
S_{\mathrm{org}}(t)&<K_{\mathrm{org}}, \nonumber\\
F(t)&\geq F_{\min}.
\label{eq:classifier}
\end{align}
The values $K_M$, $K_{\mathrm{org}}$, and $F_{\min}$ are versioned calibration
parameters. They are not universal constants.

\subsection{Replay boundary}

A complete replay parses the receipt with strict duplicate-key and number
handling, validates its schema, recomputes both digests, verifies the chain and
evidence hashes, loads the exact metric profile, reruns the classifier, and
records any disagreement. The bundled dependency-free verifier performs
structural, hash, and predecessor-link checks. Evidence retrieval and full JSON
Schema validation remain integration responsibilities.

\section{Safety constraints}

RES=RAG yields design constraints rather than a single maximization objective:
\begin{enumerate}
  \item do not optimize RES-RAG separation toward zero;
  \item distinguish retrieved evidence from generated content;
  \item retain estimator uncertainty and missing data;
  \item do not infer machine consciousness from anthropomorphic language;
  \item do not use user agreement as an Axis II stability proxy;
  \item preserve accountable human control over consequential intervention;
  \item classify unauditable state as indeterminate rather than safe; and
  \item give any client-side narrative store export, deletion, provenance,
        consent, and access controls.
\end{enumerate}

\section{Hypotheses and falsification programme}

\renewcommand{\arraystretch}{1.16}
\begin{longtable}{>{\bfseries}p{0.07\linewidth}p{0.42\linewidth}p{0.42\linewidth}}
\caption{Testable RES=RAG and CSNP hypotheses.}
\label{tab:hypotheses}\\
\toprule
ID & Hypothesis & Example falsifier\\
\midrule
\endfirsthead
\toprule
ID & Hypothesis & Example falsifier\\
\midrule
\endhead
H1 & Task quality peaks within a calibrated non-zero $W_2$ band.
   & Quality is monotonic in distance or unaffected across preregistered tasks.\\
H2 & $D_r>1$ predicts governability loss before semantic failure.
   & It has no out-of-sample predictive value over matched baselines.\\
H3 & Slowdown improves recovery under high rate pressure.
   & Randomized slowdown does not improve return-to-band time.\\
H4 & Relevant retrieval reduces above-band drift.
   & Retrieval worsens or does not change controlled grounded accuracy.\\
H5 & Diversity interventions improve below-band sterility.
   & Novelty-adjusted task quality does not improve.\\
H6 & Axis II adds predictive value beyond separate Axis I and III measures.
   & A joint model does not outperform preregistered separate baselines.\\
H7 & Hash-chained receipts improve incident reconstruction.
   & Blinded teams reconstruct incidents equally well without receipts.\\
H8 & Style-only politeness is insufficient to recover high-density,
     high-rate states.
   & Style-only intervention matches rate- or density-targeted intervention.\\
\bottomrule
\end{longtable}

Each study should preregister the representation, encoder, ground metric,
window, thresholds, missing-data rule, intervention set, outcome, and baseline.
Thresholds fitted and tested on the same data are not independent validation.

\section{Limitations}

\begin{itemize}
  \item Wasserstein values are representation-dependent; encoders and ground
        metrics can produce incompatible distances.
  \item Interface proxies do not identify internal mental or computational
        states without additional assumptions.
  \item The framework establishes neither human nor machine phenomenal
        consciousness.
  \item ``Quantum-like'' denotes contextual or non-classical behavioral
        structure only, not physical quantum computation.
  \item The dual-operator equation is phenomenological and may fail for
        discrete, path-dependent, or non-stationary systems.
  \item Deterministic receipts establish provenance integrity, not semantic
        truth.
  \item This release is a specification and verifier, not a benchmark,
        clinical instrument, or production safety certification.
\end{itemize}

\section{Research pathway}

A disciplined programme proceeds from synthetic traces to held-out studies:
\begin{enumerate}
  \item construct traces with known drift and recovery events;
  \item compare declared semantic representations and sensitivity to ground
        metrics;
  \item calibrate thresholds on development data and freeze them;
  \item test \cref{tab:hypotheses} on held-out human-machine interactions;
  \item publish receipts, metric profiles, code, uncertainty, and negative
        results; and
  \item evaluate client-controlled memory under privacy, deletion, and
        adversarial-integrity tests.
\end{enumerate}
Claims of broad societal impact should follow reproducible evidence. The
immediate contribution is narrower: RES=RAG now has a coherent mathematical
interpretation, an operational protocol, explicit falsifiers, and an auditable
implementation boundary.

\section*{Authorship and contributions}

The conceptual origin of RES=RAG, its receptive-generative distinction, and
its three-axis relational structure are attributed to Jean-Charles J. C.
Tassan. The computational interpretation, optimal-transport architecture,
history-dependent information model, operational variables, deterministic
protocol profile, validation boundary, and release engineering are attributed
to Trent Slade, developed in collaboration with Tassan.

The repository contains the proposed CRediT-compatible statement. In summary:
Tassan is proposed as lead for conceptualization and supervision, supporting
methodology, and equal for original drafting and review; Slade is proposed as
lead for methodology, formal analysis, software, validation, visualization,
and project administration, supporting conceptualization, and equal for
original drafting and review. CRediT describes contribution rather than
determining authorship \citep{brand2022credit}. Both authors must approve the
final author order, contribution statement, affiliations, and archival
metadata before DOI deposit.

\section*{Artifact availability}

The normative schema, protocol, verifier, tamper tests, Markdown manuscript,
LaTeX source, bibliography, and compiled PDF are maintained in the public
repository:
\begin{center}
\url{https://github.com/QSOLKCB/res-rag}
\end{center}

\bibliographystyle{plainnat}
\bibliography{references}

\appendix
\section{Normative artifact map}

\begin{table}[ht]
\centering
\small
\begin{tabularx}{\linewidth}{>{\raggedright\arraybackslash}p{0.48\linewidth}X}
\toprule
\textbf{Path} & \textbf{Role}\\
\midrule
\path{paper/RES_RAG_CSNP_Formalization_v1.0.0.tex}
  & Archival LaTeX source.\\
\path{paper/RES_RAG_CSNP_Formalization_v1.0.0.pdf}
  & Compiled review copy.\\
\path{paper/RES_RAG_CSNP_Formalization_v1.0.0.md}
  & GitHub-readable manuscript.\\
\path{paper/CONTRIBUTIONS.md}
  & Proposed authorship and provenance record.\\
\path{protocol/CSNP-RP_v1.0.0.md}
  & Normative protocol rules.\\
\path{protocol/csnp-rp-v1.0.0.schema.json}
  & Receipt JSON Schema.\\
\path{protocol/examples/csnp-receipt.example.json}
  & Conforming synthetic receipt.\\
\path{protocol/tools/verify_receipt.mjs}
  & Dependency-free integrity verifier.\\
\path{protocol/tests/verify_receipt.test.mjs}
  & Validity and tamper tests.\\
\bottomrule
\end{tabularx}
\caption{Version 1.0.0 artifact set.}
\end{table}

\section{Implementation note}

The canonicalization profile is intentionally small enough for independent
implementation. It is not a substitute for a general-purpose canonical JSON
standard. Cross-language deployments should first test a shared corpus of
strings, integer and fixed-precision values, arrays, and nested objects.
Changing numeric semantics or key ordering requires a new CSNP-RP version.

\end{document}
